\chapter{The Atom as a Resonant Geometric System}
\label{ch:atom-resonant}

% MERGE from Book_3/Chapter_7_Atom_And_Chemistry.tex + Book_2/Chapter_4 atomic sections

\section{Atomic Structure in SDT}

In SDT, an atom is not a probability cloud surrounding a point nucleus. It is a structured, resonant geometric system: a central nuclear vortex (the proton complex) surrounded by electron vortices locked into specific geometric configurations by the displacement field.

\subsection{The Nuclear Vortex}

The nucleus is a composite displacement vortex. Each proton is a toroidal vortex with charge radius $R_p = 0.8414$~fm. Neutrons are proton-electron composites (Chapter~\ref{ch:neutron}), geometrically locked within the nuclear structure.

The nuclear pressure field radiates outward, creating the displacement landscape in which electrons must find stable orbits.

\subsection{Electron Orbits as Geometric Resonances}

An electron is a displacement vortex locked into a resonant orbit by two competing effects:
\begin{enumerate}
  \item \textbf{Inward:} The nuclear pressure gradient (``gravity'' at the atomic scale) pulls the electron toward the nucleus.
  \item \textbf{Outward:} The electron's own displacement field creates a centrifugal pressure barrier.
\end{enumerate}

Stable orbits exist where these two pressures balance---resonant nodes in the displacement field. The velocity formula $v = (c/\kop)\sqrt{Z_{\text{eff}} \cdot R_p / r}$ gives the velocity at each resonant node.


\section{Shell Structure from Geometry}

The shell structure of atoms ($n = 1, 2, 3, \ldots$) is not an arbitrary quantum number. It reflects the geometric constraint that stable vortex orbits must satisfy:
\begin{equation}
  2\pi r_n = n \cdot \lambda_{\text{dB}}
\end{equation}

where $\lambda_{\text{dB}} = h/(m_e v_n)$ is the de Broglie wavelength. In SDT terms, this is a resonance condition: the displacement wave created by the orbiting electron must constructively interfere with itself after one complete revolution.


\section{The Electron Occupation Rules}

Each shell $n$ can hold $2n^2$ electrons:
\begin{center}
\begin{tabular}{rrl}
\toprule
$n$ & Capacity & \textbf{Subshells} \\
\midrule
1 & 2   & 1s \\
2 & 8   & 2s, 2p \\
3 & 18  & 3s, 3p, 3d \\
4 & 32  & 4s, 4p, 4d, 4f \\
\bottomrule
\end{tabular}
\end{center}

In SDT, each subshell represents a distinct geometric tiling of solid angle. The s-orbital is spherically symmetric (1 orientation). The p-orbitals tile along three axes (3 orientations). The d-orbitals tile between axes (5 orientations). The f-orbitals tile along body diagonals (7 orientations).

The factor of 2 (spin) represents the two possible helicities of the electron vortex: clockwise and counterclockwise about its axis of orbital rotation.


\section{Chemical Bonding}

Chemical bonds form when adjacent atoms' displacement fields overlap and find a lower-energy geometric configuration than the isolated atoms.

\begin{itemize}
  \item \textbf{Covalent bonds:} Two electron vortices from different atoms lock into a shared resonant orbit between the nuclei. The shared orbit reduces the total displacement energy.
  \item \textbf{Ionic bonds:} One atom's valence electron transfers entirely to another atom, achieving closed-shell (maximum geometric stability) configurations for both.
  \item \textbf{Metallic bonds:} Valence electrons delocalise across a lattice, forming a ``sea'' of mobile displacement vortices---the conduction band.
\end{itemize}


\section{Summary}

\begin{enumerate}
  \item Atoms are resonant geometric systems, not probability clouds.
  \item Shell structure reflects constructive interference of displacement waves.
  \item Subshells correspond to distinct geometric tilings of solid angle.
  \item Chemical bonding is the geometric optimisation of displacement energy.
  \item The koppa constant $\kop = 0.5464$ governs all orbital velocities identically.
\end{enumerate}
